% !TEX root = ../notes_template.tex
\chapter*{Preface}
\addcontentsline{toc}{chapter}{Preface}

The \textit{Neuromuscular \& Vascular Lab} is a foundational component of your first semester in the Plymouth State University Doctor of Physical Therapy Program. While embedded in \textbf{PTH6110 Clinical \& Functional Anatomy}, this lab serves as the anatomical counterpart to \textbf{PTH6111 Clinical Physiology: A Muscle Centered Approach}. Together, these two courses provide a uniquely integrated experience—merging structure, function, and clinical application from the very beginning of your professional training.

This manual is your lab companion, designed to help you not only locate structures but understand how they behave, interact, and adapt in clinical contexts. Just as \textit{Clinical Physiology} centers on the muscle fiber as the generative engine of movement, this lab allows you to palpate, assess, and visualize that generation in action. When you test a muscle’s performance through manual testing, evaluate segmental contributions via reflexes and pulses, or visualize tissue with rehabilitative ultrasound imaging (RUSI), you are putting physiological principles into practice.

The lab and physiology course reinforce a shared set of core principles:
\begin{itemize}
    \item \textbf{Tension} is central to both courses: generated by muscle fibers, regulated by motor units, supported by vascular and respiratory systems.
    \item \textbf{Excitation and regulation}, key themes in physiology, come alive as you trace neurologic pathways from spinal roots to peripheral nerves to muscles during manual muscle testing.
    \item \textbf{Circulation and microcirculation}, which support extracellular fluid homeostasis in physiology, are examined in the lab through pulse assessment and RUSI.
    \item \textbf{Systems thinking}, emphasized in both, is not optional but essential: movement and dysfunction must be understood through multiregional and multisystem interactions.
\end{itemize}

This lab doesn’t just teach you where things are—it challenges you to think like a clinician. You'll begin to reason through the “why” behind what you see, feel, and test. As you explore the shoulder, pelvis, or cervical spine, remember that no movement occurs in isolation. Clinical reasoning requires that you understand both the anatomy and the physiology of functional—and dysfunctional movement—as a dynamic whole.

Use this manual actively. Annotate it. Use it to generate hypotheses. Bring insights from physiology into the lab and bring questions from the lab into your physiology discussions. When you understand how a muscle creates and sustains tension—and what can compromise that process—you are beginning to think like a physical therapist.

\vspace{2em}
\noindent\textit{Sean Collins}\\
\textit{May 2025}
